%%%%%%%%%%%%%
%% MISE EN PAGE %%
%%%%%%%%%%%%%

%% Styles des sections
\newcommand{\periodafter}[1]{\ifstrempty{#1}{}{#1.}}
\titleformat{\section}[block]{\scshape\filcenter\LARGE}{\thesection.}{.5em}{}
\titleformat{\subsection}[block]{\bfseries\filcenter\large}{\thesubsection.}{.5em}{\medskip}
\titleformat{\subsubsection}[runin]{\bfseries}{\thesubsubsection.}{.5em}{\periodafter}%{}[.]
\titlespacing{\subsubsection}{0pt}{\topsep}{.5em}

%%% Styles des théorèmes
%\newtheoremstyle{ntheorem}%
%	{\topsep}{\topsep}{\itshape}{0pt}{\bfseries}{.}{.5em}%
%	{\thmnumber{#2.\hspace{.5em}}\thmname{#1}\thmnote{ (#3)}}
%	
%\newtheoremstyle{ndefinition}%
%	{\topsep}{\topsep}{\normalfont}{0pt}{\bfseries}{.}{.5em}%
%	{\thmnumber{#2.\hspace{.5em}}\thmname{#1}\thmnote{ (#3)}}
%	
%\newtheoremstyle{nremark}%
%	{\topsep}{\topsep}{\normalfont}{0pt}{\itshape}{.}{.5em}%
%	{\thmnumber{}\thmname{#1}\thmnote{ (#3)}}


%% Styles des sections
\titleformat{\section}[block]{\scshape\filcenter\Large}{\thesection.}{.5em}{}
\titleformat{\subsection}[runin]{\bfseries}{\thesubsection.}{.5em}{}[.]
\titleformat{\subsubsection}[runin]{\bfseries}{\thesubsubsection.}{.5em}{}[.]
\titlespacing{\subsubsection}{0pt}{10pt}{.5em}


%% Styles des théorèmes
%\theoremstyle{ntheorem}
\theoremstyle{plain}
  	\newtheorem{theorem}[subsubsection]{Theorem}
  	\newtheorem{proposition}[subsubsection]{Proposition}
	\newtheorem{lemma}[subsubsection]{Lemma}
  	\newtheorem{corollary}[subsubsection]{Corollary}

%\theoremstyle{ntheorem}
  %	\newtheorem{theorem}[subsubsection]{Theorem}
  %	\newtheorem{proposition}[subsubsection]{Proposition}
	%\newtheorem{lemma}[subsubsection]{Lemma}
  	%\newtheorem{corollary}[subsubsection]{Corollary}

\theoremstyle{definition}
	\newtheorem{definition}[subsubsection]{Definition}
	\newtheorem{notation}[subsubsection]{Notation}
	\newtheorem{example}[subsubsection]{Example}
	\newtheorem{remark}[subsubsection]{Remark}
	\newtheorem{Algorithm}[subsubsection]{Algorithm}
	\newtheorem{property}[subsubsection]{Property}
	
	\makeatletter
\def\@equationname{equation}
\newenvironment{eqn}[1]{%
    \def\mymathenvironmenttouse{#1}%
    \ifx\mymathenvironmenttouse\@equationname%
        \refstepcounter{subsubsection}%
    \else
        \patchcmd{\@arrayparboxrestore}{equation}{subsubsection}{}{}%          doesn't change output?
        \patchcmd{\print@eqnum}{equation}{subsubsection}{}{}%
        \patchcmd{\incr@eqnum}{equation}{subsubsection}{}{}%
%       \def\print@eqnum{\tagform@\thematrix}%                          instead of etoolbox' \pathcmd
%       \def\incr@eqnum{\refstepcounter{matrix}\let\incr@eqnum\@empty}% instead of etoolbox' \pathcmd
    \fi
    \csname\mymathenvironmenttouse\endcsname%
}{%
    \ifx\mymathenvironmenttouse\@equationname%
        \tag{\thesubsubsection}%
    \fi
    \csname end\mymathenvironmenttouse\endcsname%
}
\makeatother

%% Mise en page
\pagestyle{fancy}
%\setlength{\oddsidemargin}{0.49cm}
%\setlength{\evensidemargin}{0.49cm}
%\setlength{\topmargin}{0cm} 
%\setlength{\headheight}{1cm}
%\setlength{\headsep}{1cm}
%\setlength{\textwidth}{15cm}
%\setlength{\marginparwidth}{0cm}
%\setlength{\footskip}{2cm}
%\setlength{\headwidth}{15cm}
\setlength{\oddsidemargin}{0cm}
\setlength{\evensidemargin}{0cm}
\setlength{\topmargin}{-1.6cm}
\setlength{\headheight}{1cm}
\setlength{\headsep}{1cm}
\setlength{\textwidth}{16cm}
\setlength{\marginparwidth}{0cm}
\setlength{\footskip}{2cm}
\setlength{\headwidth}{16cm}
\newcommand{\emptysectionmark}[1]{\markboth{\textbf{#1}}{\textbf{#1}}}
\renewcommand{\sectionmark}[1]{\markboth{\textbf{\thesection. #1}}{\textbf{\thesection. #1}}}
\renewcommand{\subsectionmark}[1]{\markright{\textbf{\thesubsection. #1}}}
\fancyhead{}\fancyfoot[LC,RC]{}
\fancyhead[LE]{\leftmark}
\fancyhead[RO]{\rightmark}
\fancyfoot[LE,RO]{$\thepage$}
\fancypagestyle{plain}{
\fancyhf{}\fancyfoot[LC,RC]{}
\fancyfoot[LE,RO]{$\thepage$}
\renewcommand{\headrulewidth}{0pt}
\renewcommand{\footrulewidth}{0pt}}
\newcommand{\headrulelength}{16cm}
\renewcommand{\headrulewidth}{1pt}
\renewcommand{\footrulewidth}{0pt}
\setlength{\arraycolsep}{1pt}
\renewcommand{\headrule}{
  {\hrule width\headwidth height\headrulewidth
   \vskip-\headrulewidth}
}

%%%%%%Macro texte%%%%%%
\newcommand{\pdf}[1]{\texorpdfstring{$#1$}{1}}



\newcount\hh
\newcount\mm
\mm=\time
\hh=\time
\divide\hh by 60
\divide\mm by 60
\multiply\mm by 60
\mm=-\mm
\advance\mm by \time
\def\hhmm{\number\hh:\ifnum\mm<10{}0\fi\number\mm}

%%%%%%%%
%% XYPIC %%
%%%%%%%%

%% Options 
\renewcommand{\objectstyle}{\displaystyle}
\renewcommand{\labelstyle}{\displaystyle}
\UseTips
\SelectTips{eu}{11}

%% Nouvelle direction pour les monomorphismes (>->)
\newdir{ >}{{}*!/-10pt/@{>}}
\newdir{ -}{{}*!/-10pt/@{}}
\newdir{> }{{}*!/+10pt/@{>}}

%% Ajouts de la variante 4 pour les blancs, les lignes, les pointill�s, 
%% les zigouigouis (pas top) et les t�tes de fl�ches (presque finies).
%% Concr�tement, dans \xymatrix, la commande \ar @4 produit une 4-cellule.
\makeatletter

% Noms des nouvelles directions
\xyletcsnamecsname@{dir4{}}{dir{}}
\xydefcsname@{dir4{-}}{\line@ \quadruple@\xydashh@}
\xydefcsname@{dir4{.}}{\point@ \quadruple@\xydashh@}
\xydefcsname@{dir4{~}}{\squiggle@ \quadruple@\xybsqlh@}
\xydefcsname@{dir4{>}}{\Tttip@}
\xydefcsname@{dir4{<}}{\reverseDirection@\Tttip@}

% Commande quadruple
% Attention : la position des pourcents est importante !!
\xydef@\quadruple@#1{%
	\edef\Drop@@{%
		\dimen@=#1\relax
		\dimen@=.5\dimen@
		\A@=-\sinDirection\dimen@
		\B@=\cosDirection\dimen@
		\setboxz@h{%
			\setbox2=\hbox{\kern3\A@\raise3\B@\copy\z@}%
			\dp2=\z@ \ht2=\z@ \wd2=\z@ \box2
			\setbox2=\hbox{\kern\A@\raise\B@\copy\z@}%
			\dp2=\z@ \ht2=\z@ \wd2=\z@ \box2
			\setbox2=\hbox{\kern-\A@\raise-\B@\copy\z@}%
			\dp2=\z@ \ht2=\z@ \wd2=\z@ \box2
			\setbox2=\hbox{\kern-3\A@\raise-3\B@ \noexpand\boxz@}%
			\dp2=\z@ \ht2=\z@ \wd2=\z@ \box2
		}%
		\ht\z@=\z@ \dp\z@=\z@ \wd\z@=\z@ \noexpand\styledboxz@
	}%
}

% T�tes de quadruples fl�ches : positionnement ok, reste la jointure
% Les parties comment�es sont des essais � reprendre plus tard
\xydef@\Tttip@{\kern2pt \vrule height2pt depth2pt width\z@
	\Tttip@@ \kern2pt \egroup
	\U@c=0pt \D@c=0pt \L@c=0pt \R@c=0pt \Edge@c={\circleEdge}%
	\def\Leftness@{.5}\def\Upness@{.5}%
	\def\Drop@@{\styledboxz@}\def\Connect@@{\straight@{\dottedSpread@\jot}}}
	
\xydef@\Tttip@@{%
	\dimen@=.25\dimen@
%	\A@=-\sinDirection\dimen@
 	\B@=\cosDirection\dimen@
	\setboxz@h\bgroup\reverseDirection@\line@ \wdz@=\z@ \ht\z@=\z@ \dp\z@=\z@
%	\kern\A@ \raise\B@ \boxz@ \kern\L@c
%	\kern-\L@c \boxz@ \kern\L@c
	{\vDirection@(1,-1)\xydashl@ \xyatipfont\char\DirectionChar}%
	{\vDirection@(1,+1)\xydashl@ \xybtipfont\char\DirectionChar}%
}

% Red�finition de la commande \ar 
% Provoque un avertissement � la compilation
\xydef@\ar@form{
	\ifx \space@\next \expandafter\DN@\space{\xyFN@\ar@form}%
	\else\ifx ^\next \DN@ ^{\xyFN@\ar@style}\edef\arvariant@@{\string^}%
	\else\ifx _\next \DN@ _{\xyFN@\ar@style}\edef\arvariant@@{\string_}%
	\else\ifx 0\next \DN@ 0{\xyFN@\ar@style}\def\arvariant@@{0}%
	\else\ifx 1\next \DN@ 1{\xyFN@\ar@style}\def\arvariant@@{1}%
	\else\ifx 2\next \DN@ 2{\xyFN@\ar@style}\def\arvariant@@{2}%
	\else\ifx 3\next \DN@ 3{\xyFN@\ar@style}\def\arvariant@@{3}%
	\else\ifx 4\next \DN@ 4{\xyFN@\ar@style}\def\arvariant@@{4}%
	\else\ifx \bgroup\next \let\next@=\ar@style
	\else\ifx [\next \DN@[##1]{\ar@modifiers{[##1]}}%]
	\else\ifx *\next \DN@ *{\ar@modifiers}%
	\else\addLT@\ifx\next \let\next@=\ar@slide
	\else\ifx /\next \let\next@=\ar@curveslash
	\else\ifx (\next \let\next@=\ar@curveinout %)
	\else\addRQ@\ifx\next \addRQ@\DN@{\ar@curve@}%
	\else\addLQ@\ifx\next \addLQ@\DN@{\xyFN@\ar@curve}%
	\else\addDASH@\ifx\next \addDASH@\DN@{\defarstem@-\xyFN@\ar@}%
	\else\addEQ@\ifx\next \addEQ@\DN@{\def\arvariant@@{2}\defarstem@-\xyFN@\ar@}%
	\else\addDOT@\ifx\next \addDOT@\DN@{\defarstem@.\xyFN@\ar@}%
	\else\ifx :\next \DN@:{\def\arvariant@@{2}\defarstem@.\xyFN@\ar@}%
	\else\ifx ~\next \DN@~{\defarstem@~\xyFN@\ar@}%
	\else\ifx !\next \DN@!{\dasharstem@\xyFN@\ar@}%
	\else\ifx ?\next \DN@?{\ar@upsidedown\xyFN@\ar@}%
	\else \let\next@=\ar@error
	\fi\fi\fi\fi\fi\fi\fi\fi\fi\fi\fi\fi\fi\fi\fi\fi\fi\fi\fi\fi\fi\fi\fi \next@}

\makeatother


%%%%%%%%%%%%%%%%%
%% SYMBOLES %%%%%%%%
%%%%%%%%%%%%%%%%%

%% Flèches
\newcommand{\fl}{\rightarrow}
\newcommand{\flg}{\leftarrow}
\newcommand{\fll}{\longrightarrow}
\newcommand{\ofl}[1]{\overset{\displaystyle #1}{\fll}}
\newcommand{\ifl}{\rightarrowtail}
\newcommand{\pfl}{\twoheadrightarrow}
\newcommand{\fllg}{\longleftarrow}
\newcommand{\oflg}[1]{\overset{\displaystyle #1}{\fllg}}

\newcommand{\dfl}{\Rightarrow}
\newcommand{\dfll}{\Longrightarrow}
\newcommand{\odfl}[1]{\overset{\displaystyle #1}{\dfll}}

\newcommand{\tfl}{\Rrightarrow}
\newcommand{\qfl}{\xymatrix@1@C=10pt{\ar@4 [r] &}}

%% Parenthèses et crochets
\newcommand{\ens}[1]{\left\{ #1 \right\}}
\renewcommand{\angle}[1]{\langle #1 \rangle}
\newcommand{\pres}[2]{\langle #1 \;\vert\; #2 \rangle}
\newcommand{\bigpres}[2]{\big\langle\; #1 \;\;\big\vert\;\; #2 \;\big\rangle}
\newcommand{\cohpres}[3]{\langle #1 \;\vert\; #2 \;\vert\; #3 \rangle}
\newcommand{\bigcohpres}[3]{\big\langle\; #1 \;\;\big\vert\;\; #2 \;\;\big\vert\;\; #3 \;\big\rangle}

%% Accents
\newcommand{\op}[1]{{#1}^{o}}
\newcommand{\env}[1]{{#1}^{e}}
\newcommand{\cl}[1]{\overline{#1}}
\newcommand{\rep}[1]{\widehat{#1}}
\renewcommand{\tilde}[1]{\widetilde{#1}}
\newcommand{\tck}[1]{#1^{\top}}
\newcommand{\vect}[1]{\overrightarrow{#1}}

%% Opérateurs
\DeclareMathOperator{\id}{Id}
\DeclareMathOperator{\Ker}{Ker}
\DeclareMathOperator{\coker}{coker}
\DeclareMathOperator{\Coker}{Coker}

\DeclareMathOperator{\End}{End}
\DeclareMathOperator{\Aut}{Aut}
\DeclareMathOperator{\Lan}{Lan}

\DeclareMathOperator{\h}{h}
\DeclareMathOperator{\crit}{c}
\DeclareMathOperator{\Hom}{\mathrm{Hom}}
\DeclareMathOperator{\Ext}{\mathrm{Ext}}
\DeclareMathOperator{\Tor}{\mathrm{Tor}}
\DeclareMathOperator{\Der}{\mathrm{Der}}

\DeclareMathOperator{\FDT}{FDT}
\DeclareMathOperator{\FDTAB}{FDT_{\mathrm{ab}}}
\DeclareMathOperator{\FP}{FP}
\DeclareMathOperator{\FHT}{FHT}
\DeclareMathOperator{\U}{\mathcal{U}}
\DeclareMathOperator{\Tc}{\mathcal{T}}
\DeclareMathOperator{\Si}{\mathcal{S}}
\DeclareMathOperator{\Grw}{Grw}
\DeclareMathOperator{\Rec}{Rect}
\DeclareMathOperator{\Sh}{S}
\DeclareMathOperator{\ShST}{ShSt}
\DeclareMathOperator{\ShT}{ShT}
\DeclareMathOperator{\YoungST}{Ys}
\DeclareMathOperator{\YoungT}{Yt}
\DeclareMathOperator{\SkewT}{St}
\DeclareMathOperator{\RSK}{sRSK}
\DeclareMathOperator{\col}{Col}
\DeclareMathOperator{\Ba}{Ba}
\DeclareMathOperator{\ir}{ir}
\DeclareMathOperator{\irm}{ir_{\text{\upshape{m}}}}
\DeclareMathOperator{\lm}{lm}
\DeclareMathOperator{\lc}{lc}
\DeclareMathOperator{\lt}{lt}
\DeclareMathOperator{\wt}{wt}
\DeclareMathOperator{\up}{up}
\DeclareMathOperator{\down}{down}
\DeclareMathOperator{\sh}{sh}
\DeclareMathOperator{\std}{std}
\DeclareMathOperator{\height}{ht}
\DeclareMathOperator{\rowword}{r}
\DeclareMathOperator{\NF}{NF}
\DeclareMathOperator{\Crit}{Crit}

%% Opérations binaires
\newcommand{\tens}{\otimes}

%% Caractères 
\renewcommand{\phi}{\varphi}
\renewcommand{\epsilon}{\varepsilon}

\newcommand{\Nb}{\mathbb{N}}
\newcommand{\Zb}{\mathbb{Z}}

\newcommand{\dr}{\partial}

\newcommand{\Ar}{\EuScript{A}}
\newcommand{\Br}{\EuScript{B}}
\newcommand{\Cr}{\EuScript{C}}
\newcommand{\Dr}{\EuScript{D}}
\newcommand{\Er}{\EuScript{E}}
\newcommand{\Fr}{\EuScript{F}}
\newcommand{\Gr}{\EuScript{G}}
\newcommand{\Hr}{\EuScript{H}}
\newcommand{\Ir}{\EuScript{I}}
\newcommand{\Kr}{\EuScript{K}}
\newcommand{\Lr}{\EuScript{L}}
\newcommand{\Mr}{\EuScript{M}}
\newcommand{\Nr}{\EuScript{N}}
\newcommand{\Or}{\EuScript{O}}
\renewcommand{\Pr}{\EuScript{P}}
\newcommand{\Qr}{\EuScript{Q}}
\newcommand{\Sr}{\EuScript{S}}
\newcommand{\Rr}{\EuScript{R}}
\newcommand{\Tr}{\EuScript{T}}
\newcommand{\Ur}{\EuScript{U}}
\newcommand{\Vr}{\EuScript{V}}
%\newcommand{\Wr}{\EuScript{W}}
\newcommand{\Yr}{\EuScript{Y}}

\newcommand{\A}{\mathbf{A}}
\newcommand{\B}{\mathbf{B}}
\newcommand{\C}{\mathbf{C}}
\newcommand{\D}{\mathbf{D}}
\newcommand{\F}{\mathbf{F}}
\newcommand{\G}{\mathbf{G}}
\newcommand{\K}{\mathbb{K}}
\newcommand{\M}{\mathbf{M}}
\renewcommand{\S}{\mathbf{S}}
\newcommand{\V}{\mathbf{V}}
\newcommand{\W}{\mathbf{W}}
\newcommand{\ii}{\mathbf{i}}

%% Catégories
\def\catego#1{{\bf{\sf #1}}}

\newcommand{\Ab}{\catego{Ab}}
\newcommand{\Gp}{\mathbf{Gp}}
\newcommand{\Cat}{\catego{Cat}}
\newcommand{\Act}{\catego{Act}}
\newcommand{\Mod}{\catego{Mod}}
\newcommand{\Nat}{\catego{Nat}}
\newcommand{\Set}{\catego{Set}}
\newcommand{\Ord}{\mathbf{Ord}}
\newcommand{\Rep}{\catego{Rep}}
\newcommand{\LinCat}{\catego{LinCat}}
\newcommand{\Vect}{\catego{Vect}}
\newcommand{\Vectg}{\catego{GrVect}}
\newcommand{\Alg}{\catego{Alg}}
\newcommand{\Algg}{\catego{GrAlg}}
\newcommand{\Ch}{\catego{Ch}}
\newcommand{\Chg}{\catego{GrCh}}
\newcommand{\Grph}{\catego{Grph}}

\newcommand{\Ct}[1]{\mathrm{Ct}(#1)}
\newcommand{\Wk}[1]{\mathrm{Wk}(#1)}

\newcommand{\As}{\mathbf{As}}
\newcommand{\Mon}{\mathbf{Mon}}
\newcommand{\Perm}{\mathbf{Perm}}
\newcommand{\Brd}{\mathbf{Braid}}
\newcommand{\Colo}{\mathbf{Coln}}
\newcommand{\Std}{\mathbf{Std}}
\newcommand{\inv}{\operatorname{inv}}
\newcommand{\Inv}{\operatorname{Inv}}

%% Cohomologie
\newcommand{\fnat}[2]{F_{#1}[#2]}
\newcommand{\ho}{\mathrm{H}}
\newcommand{\poldim}[1]{d_{\mathrm{pol}}(#1)}
\newcommand{\cohdim}[1]{\mathrm{cd}(#1)}
%


\renewcommand{\objectstyle}{\displaystyle}
\renewcommand{\labelstyle}{\displaystyle}
\renewcommand{\labelitemi}{$-$}
\renewcommand{\theenumi}{\roman{enumi}}
\renewcommand{\labelenumi}{\textup{\textbf{\theenumi)}}}
%\renewcommand{\labelenumi}{\upshape{\Textbf{\theenumi)}}}
\newcommand{\tableau}[1]{\texttt{#1}}
\newcommand{\Cdot}{\!\cdot\!}
\renewcommand{\leq}{\leqslant}
\renewcommand{\geq}{\geqslant}
\renewcommand{\rep}[1]{\widehat{#1}}
\DeclareMathOperator{\rem}{Rem}
\newcommand{\Pl}{\mathbf{P}}
\newcommand{\HP}{\mathbf{Hp}}
\newcommand{\insl}[1]{\rightsquigarrow_{#1}}
\newcommand{\insr}[1]{\;\raisebox{0.1em}{{\rotatebox[origin=c]{180}{$\rightsquigarrow$}}}_{#1}\;}
\newcommand{\Right}[1]{\ensuremath{\xrightarrow{ \ #1\ }}}
\newcommand{\Left}[1]{\ensuremath{\xleftarrow{ \ #1\ }}}
\newcommand{\Up}[1]{\ensuremath{\big \uparrow \!\text{\raisebox{.1ex}{\scriptsize $#1$}}}}
\newcommand{\Down}[1]{\ensuremath{\big \downarrow \!\text{\raisebox{.1ex}{\scriptsize $#1$}}}}
\newcommand{\DRight}[1]{\ensuremath{\xdashrightarrow{ \ #1\ }}}

%%%%%%%%%
%%%Colors%%%
%%%%%%%%%

\definecolor{cyan}{RGB}{175,238,238} 
\definecolor{Red}{rgb}{0.96,0.17,0.20}
\definecolor{RedD}{rgb}{0.57, 0.0, 0.04}
\definecolor{Green}{rgb}{0.0,1.0,0.0}
\definecolor{GreenL}{rgb}{0.0,1.0,0.0} 
\definecolor{GreenD}{rgb}{0.64,0.76,0.68} 
\definecolor{Yellow}{rgb}{1.0,1.0,0.19}
\definecolor{YellowD}{rgb}{0.93, 0.84, 0.25}
\definecolor{Blue}{rgb}{0.0,1.0,1.0}
\definecolor{BlueD}{rgb}{0.63,0.79,0.95}
\definecolor{vert}{rgb}{0,0.45,0}
\definecolor{bazaar}{rgb}{0.6, 0.47, 0.48}
\definecolor{bronze}{rgb}{0.8, 0.5, 0.2}
\definecolor{darkspringgreen}{rgb}{0.09, 0.45, 0.27}


%%%%%%%%%%%%%%%%%%%%%
\def\len{\ell}
\def\Hr{H_r}
\def\Hl{H_l}
\def\Ich{C_r}
\def\Jch{C_l}
\def\Ar{\mathcal{A}}
\def\Br{\mathcal{B}}
\def\Sr{\mathcal{S}}
\def\Fr{\mathcal{F}}
\def\Rr{\mathcal{R}}
\def\Lr{\mathcal{L}}
\def\Pr{\mathcal{P}}
\def\Ir{\mathcal{I}}
\def\Tr{\mathcal{T}}
\def\Yr{\mathcal{Y}}
\makeatletter
\def\blfootnote{\xdef\@thefnmark{}\@footnotetext}
\makeatother
\def\R{\mathbf{R}}
\def\P{\mathbf{P}}
\def\T{\mathbf{T}}
\def\Q{\mathbf{Q}}
\def\Sk{\mathbf{Sk}}
\def\Xr{\mathcal{X}}
\def\St{\mathbf{Std}}

\def\points{\rotatebox{90}{\hspace{-1pt}...}}
